\section{Discussion}
\label{sec:discussion}

\paragraph{The bridge: does the black-box explanation match the white-box one?} The study's
distinctive move is to ask one question two ways---which \emph{instruction} matters (Stage~1) and
which \emph{direction} matters (Stage~2)---on the same model, and then to check whether the answers
agree. That check is the correspondence between the components RQ1 finds most necessary and the
sub-directions RQ8 finds separable.
\begin{resultbranch}{bridge holds}
If RQ1 ranks C5 (negatives) and C1 (color) most necessary \emph{and} RQ8 finds those the sub-vectors
that most cleanly steer their own metric families, the two explanations converge: the parts of the
prose that carry the causal weight are the parts with the crispest internal directions. That
convergence is the paper's most satisfying result---a prompt-level account and a
representation-level account of the \emph{same} phenomenon, mutually corroborating---and it suggests
skill authors are, in effect, writing instructions that map onto the model's existing axes of
variation.
\end{resultbranch}
\begin{resultbranch}{bridge fails}
If the most necessary components are \emph{not} the most separable directions (RQ8 absent or
mismatched), the two levels tell different stories: the instruction's compositional structure is not
the model's. This is equally interesting---it says the componentization that helps a human author a
skill is not how the model represents design---and it cautions against reading prompt structure as
mechanism.
\end{resultbranch}

\paragraph{What ``make it beautiful'' tells us about instruction compression.} RQ4 turns a cheap
one-liner into a probe of how much of a 391-token skill is \emph{knowledge} versus \emph{permission}.
If a five-word nudge recovers most of the effect (RQ4 nudge), the model already contains the
capability and the skill largely licenses its use---which is precisely the premise that makes Stage~2
plausible: a capability the model already has is one a direction can actuate. If instead the
structured content is load-bearing (RQ4 content), the skill transfers specific
knowledge---named palettes, type scales, prohibitions---and we should expect the extracted direction
to be correspondingly richer than a single ``try harder'' axis. Either way, RQ4 and Stage~2 are two
readings of the same question: is the design skill flipping a switch the model already has, or
installing something new?

\paragraph{Mechanistic implications.} A positive Stage~2 result would extend the single-direction
picture---clean for refusal \citep{arditi2024refusal}---to a graded, compositional, generation-spanning
\emph{aesthetic} behavior, which is a substantially harder case than a binary safety behavior: the
target is downstream of thousands of tokens of code and has no single ``answer token.'' A null would
mark a boundary of that picture, and the honest reading is that not every behavior a prompt induces is
a steerable feature. The correspondence result (RQ8), whichever way it falls, speaks to the geometry of
compositional concepts \citep{park2024geometry}: whether a human's five-way decomposition of ``design''
lines up with the model's internal near-orthogonal directions.

\paragraph{Relation to the steering-reliability literature.} We inherit its central warning: mean
effects hide heavy tails, and up to half of inputs can be anti-steerable
\citep{tan2024analysing,dasilva2025steeringoff}. Our protocol answers it directly---per-prompt
distributions, an anti-steerable fraction, a class-separation pre-screen, an OOD held-out over unseen
\emph{types}, and a coherence guardrail on a non-monotone dose-response
\citep{taimeskhanov2026strength}. The four-conjunct success criterion, in particular, is designed so
that a strong average cannot by itself buy a causal claim: specificity and necessity must also hold.
This is the methodological transfer we hope outlasts the particular numbers---a template for claiming a
\emph{causal} steering result responsibly when the underlying phenomenon is known to be unreliable.

\paragraph{Practical guidance for skill design.} Read together, the Stage-1 outcomes yield actionable
advice for authors of design skills: invest tokens where necessity concentrates (the RQ1 ranking),
prune any component whose leave-one-out is null or negative, and exploit the interactions RQ3
surfaces---if C1$\times$C5 is super-additive, a palette instruction and a ``no purple gradients''
prohibition should ship together, not separately. If sufficiency is distributed (RQ2), resist the urge
to shorten the skill to its single largest component; if it is concentrated, the skill can be
compressed. These are empirical recommendations, contingent on which branch obtains, and each is
traceable to a preregistered contrast.
